% \begin{itemize}
%     \item scheme A/B/C
%     \item MHOU
%     \item covariance matrix
% \end{itemize}
We now attempt to address a problem we first see in \cref{sec:resum}: Missing Higher Order Uncertainty (MHOU).
We repeat again, the truncation of the perturbative series is a fundamental part of our framework and we are bound to it.
The naive solution of just computing the next higher correction is not a solution, as we would face the problem just again there.
Thus the task is to get an estimate of the missing terms based on current knowledge, which by definition is an ill-posed problem.
Several possible strategies have been suggested in the literature~\cite{NNPDF:2024dpb,Lim:2024nsk,Tackmann:2024kci,Kassabov:2022orn,Bonvini:2020xeo,Cacciari:2011ze,Duhr:2021mfd,Ghosh:2022lrf}, which are based on resummation of specific variables and/or statistical arguments.

From the practical point of view as a PDF fitter, however, there is one big problem with many of the proposed algorithms: they only work for some specific set of observables and a detailed analysis of the kinematics.
Thus, we review here the most common approach to estimate MHOU, which can be rigorously derived in pQCD and applied to any observable: scale variations.
On the other side, especially in recent years, scale variations have been challenged as suitable tool, due to the inherent arbitrariness and the lack of statistical interpretation (to which we return in the end).
For now we focus on their practical implementation in DIS.

\subsection{Renormalization scale}
We start again from the simpler case of the renormalization scale $\mu_R$ associated with the strong coupling $\alpha_s(\mu_R^2)$ via the RGE, \cref{eq:LObeta}.
As we say above, the renormalization scale captures the ultra-violet divergencies, which arise, e.g., from self-energy diagrams and which are thus associated with the partonic diagram itself.
Moreover, as we see from \cref{eq:CpQCD}, we expand the partonic coefficient functions in powers of the strong coupling.
\begin{conclusionbox}{Renormalization scale dependence}
    The renormalization scale dependence is captured by the partonic coefficient function $\vb C(z,Q^2,\mu_R^2)$.
\end{conclusionbox}
Since we have to choose something for the arbitrary scale $\mu_R$ and we only have a single scale available, the virtuality $Q^2$, we start by setting them equal to each other $\mu_R^2 = Q^2$, which is also the implicit choice in \cref{eq:fact1}.
We then need to find a way to estimate the missing terms in $\vb C(z,Q^2,\mu_R^2=Q^2)$, i.e.\ something that is identical at the requested perturbative accuracy, but differs from it beyond.
This can be achieved by using \cref{eq:LOalphasexp} together with a suitable redefinition of the perturbative coefficients of the partonic coefficient function $\vb C^{(j)}(z)$.
Specifically, we require
\begin{equation}
    \vb C(z,Q^2,\mu_R^2) = \vb C(z,Q^2,\mu_R^2=Q^2)\qty(1 + \order{\alpha_s(Q^2)})
\end{equation}
with
\begin{equation}
    \vb C(z,Q^2,\mu_R^2) = \sum_{j=0} \qty(\alpha_s(\mu_R^2))^j \vb C^{(j)}(z,\ln(\mu_R^2/Q^2)) \,. \label{eq:CbarRpQCD}
\end{equation}
Before we continue with deriving explicit expressions for $\vb C^{(j)}(z,\ln(\mu_R^2/Q^2))$ some comments are in order:
\begin{itemize}
    \item $\vb C(z,Q^2)$ in \cref{eq:fact1} is an abbreviation of $\vb C(z,Q^2,\mu_R^2=Q^2)$ and $\vb C^{(j)}(z)$ of $\vb C^{(j)}(z,0)$
    \item in \cref{eq:CbarRpQCD} $\mu_R$ has its intended meaning: it is the scale of the strong coupling
    \item from \cref{sec:as} we know $\alpha_s(\mu_R^2)$ is a pure resummation of logarithms, thus $\vb C^{(j)}$ may only depend on $\ln(\mu_R^2/Q^2)$
    \item $\vb C^{(j)}(z,0)$ can be computed from diagrams by assuming explicitly $\mu_R = Q$ there
\end{itemize}

We now derive the explicit next-to-next-to-leading order (NNLO) expression for fully inclusive DIS.
In particular, we recall that LO is given by $\vb C^{(0)}(z)$, i.e.\ proportional to $\qty(\alpha_s(Q^2))^0$ in \cref{eq:CpQCD}, NLO is given by $\vb C^{(1)}(z)$ proportional to $\qty(\alpha_s(Q^2))^1$, and NNLO by $\vb C^{(1)}(z)$ proportional to $\qty(\alpha_s(Q^2))^2$.
Moreover, we assume that the coefficient functions at the central scale $\mu_R=Q$ are known and we can write \cref{eq:CpQCD} up to NNLO
\begin{equation}
    \vb C(z,Q^2) = \vb C^{(0)}(z) + \alpha_s(Q^2) \vb C^{(1)}(z) + \qty(\alpha_s(Q^2))^2 \vb C^{(2)}(z) \label{eq:CNNLO}
\end{equation}
Expanding \cref{eq:CbarRpQCD} up to the same accuracy we get
\begin{align}
    \vb C(z,Q^2,\mu_R^2) &= \vb C^{(0)}(z,\ln(\mu_R^2/Q^2)) + \alpha_s(\mu_R^2) \vb C^{(1)}(z,\ln(\mu_R^2/Q^2)) \nonumber\\
    &\hspace{20pt} + \qty(\alpha_s(\mu_R^2))^2 \vb C^{(2)}(z,\ln(\mu_R^2/Q^2)) \label{eq:CbarRNNLO}
\end{align}
We are now left with the task to make \cref{eq:CNNLO} and \cref{eq:CbarRNNLO} perturbatively equivalent, and although we are considering NNLO the equivalence must hold at any lower order.
Thus, we first can consider LO: since there is no dependency on the strong coupling in the first place, the scale varied coefficient function may not depend on the renormalization scale and we find
\begin{equation}
    \vb C^{(0)}(z,\ln(\mu_R^2/Q^2)) = \vb C^{(0)}(z)\,.
\end{equation}
Starting from NLO an explicit dependency on the strong coupling appears, however, from \cref{eq:LOalphasexp} we know that the two choices, $\alpha_s(Q^2)$ and $\alpha_s(\mu_R^2)$ differ by higher order terms.
Specifically, we have
\begin{equation}
    \alpha_s(Q^2) = \alpha_s(\mu_R^2) + \order{\qty(\alpha_s(\mu_R^2))^2} \qq{and} \alpha_s(\mu_R^2) = \alpha_s(Q^2) + \order{\qty(\alpha_s(Q^2))^2}
\end{equation}
Thus, if \cref{eq:CNNLO} and \cref{eq:CbarRNNLO} are the same up to $\order{\qty(\alpha_s(Q^2))^2}$ we find
\begin{equation}
    \vb C^{(1)}(z,\ln(\mu_R^2/Q^2)) = \vb C^{(1)}(z)\,.
\end{equation}

Note that with this identification \cref{eq:CNNLO} and \cref{eq:CbarRNNLO} are not the same in practice, they are only perturbative equivalent.
In any actual calculation with $\mu_R\neq Q$ also $\alpha_s(\mu_R^2)$ and $\alpha_s(Q^2)$ are different: they differ by an amount which is beyond our fixed-order framework and if that difference is \enquote{small} depends on the actual values of the strong coupling of the respective scales.
However, with that we achieve the main purpose of this game: we generate some higher order terms, which we can use to estimate MHOUs.

Starting from NNLO we see that the matching becomes non-trivial.
In the case at hand we can insert \cref{eq:LOalphasexp} into \cref{eq:CNNLO} (with the choice $(\mu_R,\mu_{R,0})\to (Q,\mu_R)$) and re-expand in powers of $\alpha_s(\mu_R^2)$ as requested.
We find
\begin{equation}
    \vb C^{(2)}(z,\ln(\mu_R^2/Q^2)) = \vb C^{(2)}(z) -  \beta_0 \ln(\mu_R^2/Q^2) \vb C^{(1)}(z)\,,
\end{equation}
so the scale-varied coefficient receives a correction from a lower-order coefficient function.

Let's take a step back and review the general features.
\begin{conclusionbox}{}
    Scale-varied coefficient functions
    \begin{itemize}
        \item can always be computed a posteriori from their non-scale-varied counterparts
        \item are a linear combination of their lower order non-scale-varied counterparts
        \item are derived by re-expanding the strong coupling in the perturbative sum up to the required perturbative accuracy
    \end{itemize}
\end{conclusionbox}
We also know that the dependency on the scale variation $\ln(\mu_R^2/Q^2)$ must be a polynomial dependence since the running of strong coupling is a pure resummation of this logarithm.
By re-expanding we mean a Taylor expansion of the strong coupling at the default scale, here $\alpha_s(Q)$, in terms of the value at the renormalization scale, $\alpha_s(\mu_R^2)$.
For computing the strong coupling at any scale $\alpha_s(\mu_R^2)$ we must use required perturbative accuracy, but for deriving the scale varied coefficient functions we only need one order less.
In the special case of fully inclusive DIS, which has no dependence on the strong coupling at LO, we even need only two orders lower.

For the sake of correcting \cref{eq:fact3}, we can improve it by writing
\begin{equation}
    \sigma(x,Q^2,\mu_R^2) = \qty(\vb C(Q^2,\mu_R^2) \otimes \vb f(Q^2))(x) \,.
\end{equation}

\subsection{Factorization scale}
As in \cref{sec:DGLAP} the general algorithm repeats for the factorization scale $\mu_F$, which is the intrinsic scale of PDFs, except everything is more complicated.
However, it is not only more complicated because PDFs are a set of function, but also because the PDF explicitly appears in \cref{eq:fact1}.
This implies, that the factorization scale dependence may not only be captured by the coefficient function, in contrast to the renormalization scale.
In particular, the factorization scale captures collinear divergencies, which arise in diagrams, and re-attributes them to the PDF, thus there is a direct connection between the two.
Recall our main goal here: we want to estimate MHOU by generating higher order terms and we simply have more freedom here to do so.
In order to simplify the discussion we explicitly neglect renormalization scale variations for the moment.

The most common way to deal with factorization scale variation is to directly modify the coefficient functions, as in the case for \cref{eq:CbarRpQCD}, and we refer to this choice as \enquote{factorization scale variation scheme C}~\cite{NNPDF:2024dpb}.
Specifically, we make a similar ansatz and write
\begin{equation}
    \sigma(x,Q^2,\mu_F^2) = \qty(\vb C(Q^2,\mu_F^2) \otimes \vb f(\mu_F^2))(x)
\end{equation}
where now $\mu_F$ has its intended place and the usual perturbative expansion of the coefficient function still holds, i.e.\
\begin{equation}
    \vb C(z,Q^2,\mu_F^2) = \sum_{j=0} \qty(\alpha_s(Q^2))^j \vb C^{(j)}(z,\ln(\mu_F^2/Q^2)) \,. \label{eq:CbarFpQCD}
\end{equation}
